% chuong_v3/fig_ch08_octave_cost.tex
% Minh họa cơ chế chi phí quãng tám (Octave Cost) của Praat:
% Ưu tiên tần số cao, phạt nhẹ chu kỳ dài để phá thế cân bằng giữa tau0 và 2tau0
\begin{tikzpicture}[x=0.06cm, y=5.8cm, >=Stealth]
    % Lưới nền mờ
    \draw[very thin, gray!20, step=1] (20,0.60) grid[xstep=30, ystep=0.1] (225,1.05);

    % Trục tọa độ
    \draw[->, thick, black!70] (20,0.60) -- (235,0.60) node[right, font=\footnotesize] {Độ trễ $\tau$ (mẫu)};
    \draw[->, thick, black!70] (20,0.60) -- (20,1.08) node[above, font=\footnotesize] {Điểm tương quan / Độ mạnh};

    % Vạch chia trục hoành (phân bố đều, không chen chúc)
    \foreach \x in {26, 50, 80, 110, 160, 200, 214} {
        \draw[thick, black!60] (\x,0.61) -- (\x,0.59) node[below, font=\scriptsize, fill=white, inner sep=1.5pt] {\x};
    }
    % Vạch chia trục tung
    \foreach \y in {0.6, 0.7, 0.8, 0.9, 1.0} {
        \draw[thick, black!60] (22,\y) -- (18,\y) node[left, font=\scriptsize] {\y};
    }

    % Đường cong U[tau]: Đỉnh tương quan chuẩn hóa (chưa trừ chi phí quãng tám)
    % Tín hiệu có tau0 = 80 (f0 = 200 Hz), đỉnh tại tau0=80 (0.920) và 2tau0=160 (0.925 do nhiễu trầm)
    \draw[thick, blue!80!black, smooth] plot coordinates {
        (30,0.62) (50,0.65) (70,0.80) (78,0.915) (80,0.920) (82,0.915) (90,0.75) 
        (110,0.63) (130,0.66) (150,0.81) (158,0.920) (160,0.925) (162,0.920) (170,0.76) 
        (190,0.64) (214,0.62)
    };
    \node[blue!85!black, font=\scriptsize, fill=white, inner sep=2pt, rounded corners=1pt] at (68,0.99) 
        {Đỉnh gốc $U[80] = 0{,}920$};
    \node[blue!85!black, font=\scriptsize, fill=white, inner sep=2pt, rounded corners=1pt] at (172,0.99) 
        {Đỉnh bội $U[160] = 0{,}925$};

    % Đường phạt chi phí quãng tám: Sc[tau] = U[tau] - K[tau]
    \draw[very thick, red!85!black, dashed, smooth] plot coordinates {
        (30,0.618) (50,0.640) (70,0.786) (78,0.900) (80,0.904) (82,0.899) (90,0.732) 
        (110,0.609) (130,0.637) (150,0.785) (158,0.894) (160,0.899) (162,0.894) (170,0.733) 
        (190,0.611) (214,0.590)
    };
    % Đặt nhãn Sc[80] sang trái đỉnh tau=80 và Sc[160] sang phải đỉnh tau=160 (hoàn toàn tách biệt)
    \node[left, red!85!black, font=\scriptsize, fill=white, inner sep=2pt, rounded corners=1pt, draw=red!30] at (72,0.88) 
        {$S_c[80] = 0{,}904$};
    \node[right, red!85!black, font=\scriptsize, fill=white, inner sep=2pt, rounded corners=1pt, draw=red!30] at (168,0.88) 
        {$S_c[160] = 0{,}899$};

    % Mũi tên và nhãn phạt Octave Cost - Đặt nhãn vào khoảng không thoáng (y=0.83), cung dẫn uốn cong hoàn toàn trên đường cong
    \draw[->, very thick, purple!75!black] (80,0.920) -- (80,0.905);
    \node[purple!80!black, font=\scriptsize, fill=white, inner sep=2pt, rounded corners=1pt, draw=purple!40] (lbl_cost80) at (96,0.83) 
        {Trừ $0{,}0158$};
    \draw[->, thick, purple!75!black] (lbl_cost80.north west) to[out=135, in=-20] (81.5,0.912);

    \draw[->, very thick, purple!75!black] (160,0.925) -- (160,0.900);
    \node[purple!80!black, font=\scriptsize, fill=white, inner sep=2pt, rounded corners=1pt, draw=purple!40] (lbl_cost160) at (144,0.83) 
        {Trừ $0{,}0259$};
    \draw[->, thick, purple!75!black] (lbl_cost160.north east) to[out=45, in=200] (158.5,0.912);

    % Hộp phân định kết luận - Đặt hoàn toàn ngoài vùng đồ thị, cách xa trục hoành
    \node[draw=green!60!black, fill=green!5, rounded corners=3pt, font=\scriptsize, align=left, inner sep=6pt] at (125,0.36) {
        \textbf{Phân định người chiến thắng qua cơ chế Octave Cost:}\\[2pt]
        $\bullet$ Trước Octave Cost: $U[160] = 0{,}925 > U[80] = 0{,}920 \implies$ Thuật toán bị đánh lừa chọn nhầm nửa tần số ($100\text{ Hz}$).\\
        $\bullet$ Sau Octave Cost: $S_c[80] = 0{,}904 > S_c[160] = 0{,}899 \implies$ \textbf{Ưu tiên tần số cao, chọn đúng $\tau_0 = 80$ ($200\text{ Hz}$)}.
    };
\end{tikzpicture}
